\section{The signed radial current of the physical angular set}
\label{sec:v70-radial-current}
The loss in \eqref{eq:v70-angular-loss} can be bounded by the
negative variation of the actual angular occupation. For circular
scatterers this variation exists at every fixed word on every compact
radial interval of the selected chart. We prove this fact and identify
its signed geometric source. No count-uniform variation estimate is
inferred from finite-word definability.

Fix a selected chart, an exact label $\lambda=(n,k)$ and
$0<S<h_\chi$. In the cylinder $(0,S)\times\mathbb S^1$ set
\[
 E_{z,\lambda}=
 \{(u,\phi):B_{z,\lambda}(\sqrt{2u}\,\omega_\phi)=1\},
 \qquad \pi(u,\phi)=u.
\]
All physical and section decisions are made before this set is formed.
We take distributional derivatives in the open cylinder; no artificial
boundary mass is inserted at $u=0,S$ or at an angular coordinate cut.

\begin{theorem}[Finite-word physical boundary current]
\label{thm:v70-physical-current}
The indicator $\one_{E_{z,\lambda}}$ has bounded variation on the
cylinder, and $L_z\in BV(0,S)$. As finite signed Radon measures,
\begin{equation}\label{eq:v70-signed-current}
 DL_z=\frac1{2\pi}\pi_\#D_u\one_{E_{z,\lambda}}
     =-\frac1{2\pi}\pi_\#
       \bigl(n_u\,\mathcal H^1|_{\partial^*E_{z,\lambda}}\bigr).
\end{equation}
Here $n$ is the measure-theoretic outward normal. In particular
\begin{equation}\label{eq:v70-negative-current}
 (DL_z)^-\le\frac1{2\pi}\pi_\#(D_u\one_{E_{z,\lambda}})^-.
\end{equation}
Radial jumps and tangencies are included. There is no division by a
boundary gradient. Finiteness is asserted for each fixed word and
label, not uniformly after taking the collision-count limit.
\end{theorem}
\begin{proof}
On a closed annulus inside the Morse disk the selected contact graph
and the fixed Morse map are real analytic. Indeed the selected
incidences remain positive and the internal contact derivative remains
invertible. The analytic implicit function theorem applies at every
point, and uniqueness identifies its local solutions. The explicit
matrix-square-root Morse map is analytic on the positive definite
cone. Compactness permits a finite analytic coordinate cover.

On this cover the original first-hit decisions have a finite Boolean
description by restricted analytic tests. Competing disks are kept
separately. If the closest point of a segment changes, divide into its
interior-foot and two endpoint cases before writing the distance
inequality. Selected flight lengths stay positive. Section membership
is tested in the original finitely many analytic rectangles. Occupation,
terminal membership and the fixed exact label are finite Boolean
combinations of these same tests. Passing to angular coordinates on a
finite cover preserves restricted analyticity. Thus $E_{z,\lambda}$
is definable in the restricted analytic structure on this compact
annulus.

By cell decomposition, its angular fibers are finite unions of
intervals with definable endpoints, with finitely many changes of
radial cell. Consequently their lengths divided by $2\pi$ define a
bounded definable function of $u$. The monotonicity theorem gives
finitely many monotone pieces, so this function has finite variation.
These uses of cell decomposition and monotonicity are the standard
ones in \cite{Coste}. They give no bound independent of the defining
word. The boundary curves also have finite length: divide each
bounded definable curve into finitely many arcs on which both
coordinates are monotone; its length is bounded by the sum of their
coordinate variations. This proves finite perimeter on the annulus.

Near $u=0$, Lemma~\ref{lem:v68-root-order} gives finitely many
analytic angular boundary graphs $\phi=\phi_j(r)$ with a fixed
cyclic order, where $u=r^2/2$. Their lengths are finite, since both
$r\mapsto r^2/2$ and the analytic $\phi_j$ have integrable derivatives
on the small closed interval. The angular fraction is analytic in $r$
by Theorem~\ref{thm:v68-angular-germ-classification}, and therefore
has finite variation after this monotone change of variable.
Joining this neighborhood to one compact annulus proves the two
bounded-variation assertions on $(0,S)$. Values on boundary curves
may be changed without changing the source or the distributions.

For $\eta\in C_c^1(0,S)$, Fubini gives
\[
 -\int_0^S L_z(u)\eta'(u)\,du
 =-\frac1{2\pi}\int_{E_{z,\lambda}}\eta'(u)\,du\,d\phi
 =\frac1{2\pi}\int\eta\circ\pi\,dD_u\one_{E_{z,\lambda}}.
\]
This proves the first equality in \eqref{eq:v70-signed-current}.
The finite-perimeter derivative formula, with the stated outward
normal convention, gives the second; see \cite{AFP}. For a signed
measure $\mu=\mu^+-\mu^-$, the negative part of $\pi_\#\mu$
is at most $\pi_\#\mu^-$. Applying this elementary measure inequality
proves \eqref{eq:v70-negative-current}. No regular-level parametrization
or transversality divisor was used.
\end{proof}

\paragraph{Attribution to primitive physical decisions.}
There is a finite ordered list of the original scalar tests at each
word. Away from their zero sets all signs, hence the actual Boolean
indicator, are locally constant. The reduced boundary in
\eqref{eq:v70-signed-current} is therefore contained, up to
$\mathcal H^1$-null sets, in the union of the genuine decision
interfaces. Identically vanishing tests have no interface and cause
no extra current. Assign each boundary point to the first interface
in this list which contains it. This gives a disjoint measurable
partition of the reduced boundary, including coincident interfaces.
The current assigned there is the jump of the actual one-hot indicator,
not a sum of jumps of all tests vanishing at that point. In particular
births and losses are combined with their signs before taking a
negative part after radial projection. The original first-clearance
witness and its next-collision mark remain in $b_z$; this additional
boundary attribution does not replace either of them.

\begin{proposition}[One-sided radial loss formula]
\label{prop:v70-loss-current}
Let $2H<h_\chi$ and choose $S$ with $2H<S<h_\chi$. For almost
every $0<u<H$,
\begin{equation}\label{eq:v70-loss-current}
 \Delta_{z,(H,2H)}(u)
 \le\int_{(u,2H)}\kappa_H(s)\,d(DL_z)^-(s),
 \qquad
 \kappa_H(s)=
 \begin{cases}
 1,&0<s\le H,\\
 (2H-s)/H,&H<s<2H.
 \end{cases}
\end{equation}
The right side retains atomic radial deaths. In particular
$DL_z\ge0$ on $(0,2H)$ implies that the complete inner source
satisfies Corollary~\ref{cor:v70-no-loss-height}.
\end{proposition}
\begin{proof}
Use the precise representative of the BV function. At its continuity
points $u<v$ the fundamental theorem for BV gives
$L_z(u)-L_z(v)\le (DL_z)^-((u,v))$. The exceptional jump points
are countable and do not affect the almost-everywhere assertion or
the average in $v$. Average over $H<v<2H$ and apply Tonelli to the
nonnegative variation measure. A point $s\le H$ is crossed by the
entire averaging interval; a point $H<s<2H$ is crossed by a fraction
$(2H-s)/H$. This proves \eqref{eq:v70-loss-current}. A BV function
with nonnegative distributional derivative has a nondecreasing
representative, giving the last assertion.
\end{proof}

For example, a downward jump of size $\beta$ at $s_0\in(H,2H)$
contributes $\beta(2H-s_0)/H$ to the right side for $u<H$.
Replacing the distributional derivative by its almost-everywhere
classical derivative would lose this contribution. Conversely, upward
jumps incur no negative-current cost. This is why absolute boundary
variation is not used in \eqref{eq:v70-loss-current}.

Combining \eqref{eq:v70-loss-density} and
\eqref{eq:v70-loss-current} identifies a sufficient physically
weighted estimate for the remaining inner loss. Specifically its
height is at most
\[
 e^{K_\chi\sqrt{2H}}W_\delta(H)
 \limsup_m\sup_{R,\lambda,t}m^2
 \sum_{z:\,t_z\in[t-H,t]}J_z
            \int_{(0,2H)}\kappa_H(s)\,d(DL_z)^-(s).
\]
This expression may be infinite. Fixed-word finite perimeter proves
that its summands are defined; it does not establish the ordered
weighted bound. The exact loss \eqref{eq:v70-angular-loss}, rather
than this sufficient upper bound, is used in the full source identity
below.
